\documentclass[10pt,a4paper]{article}

\usepackage[margin=14mm,top=14mm,bottom=16mm]{geometry}
\usepackage{fontspec}
\usepackage[table]{xcolor}
\usepackage{tabularx}
\usepackage{array}
\usepackage{enumitem}
\usepackage{titlesec}
\usepackage{fancyhdr}
\usepackage{lastpage}
\usepackage{ragged2e}
\usepackage[most]{tcolorbox}
\usepackage[hidelinks]{hyperref}

\IfFontExistsTF{Arial}{
  \setmainfont{Arial}
  \setsansfont{Arial}
}{
  \setmainfont{TeX Gyre Heros}
  \setsansfont{TeX Gyre Heros}
}

\definecolor{AIEPRed}{HTML}{D62027}
\definecolor{AIEPNavy}{HTML}{102A43}
\definecolor{AIEPDeep}{HTML}{082B5C}
\definecolor{AIEPInk}{HTML}{243B53}
\definecolor{AIEPSlate}{HTML}{52606D}
\definecolor{AIEPPaper}{HTML}{F8F3EC}
\definecolor{AIEPSoft}{HTML}{EDE6DA}
\definecolor{AIEPSoftBlue}{HTML}{E6EEF7}
\definecolor{AIEPGreen}{HTML}{3FAE6A}
\definecolor{AIEPGreenSoft}{HTML}{E9F6EE}
\definecolor{AIEPGold}{HTML}{E0BC5A}
\definecolor{Line}{HTML}{D8CFC4}
\definecolor{White}{HTML}{FFFFFF}

\pagestyle{fancy}
\fancyhf{}
\renewcommand{\headrulewidth}{0pt}
\fancyfoot[L]{\scriptsize\textcolor{AIEPSlate}{Planificación docente - Taller 2}}
\fancyfoot[R]{\scriptsize\textcolor{AIEPSlate}{GeoGreen Escolar Osorno - \thepage/\pageref{LastPage}}}

\setlength{\parindent}{0pt}
\setlength{\parskip}{3.8pt}
\setlist[itemize]{leftmargin=4.2mm,topsep=1pt,itemsep=1.2pt,parsep=0pt}
\setlist[enumerate]{leftmargin=5mm,topsep=1pt,itemsep=1.2pt,parsep=0pt}
\renewcommand{\arraystretch}{1.18}
\titleformat{\section}{\Large\bfseries\color{AIEPNavy}}{}{0pt}{}
\titleformat{\subsection}{\normalsize\bfseries\color{AIEPInk}}{}{0pt}{}

\newcolumntype{Y}{>{\RaggedRight\arraybackslash}X}
\newcolumntype{L}[1]{>{\RaggedRight\arraybackslash}p{#1}}

\newtcolorbox{card}[1][]{
  enhanced,
  colback=White,
  colframe=Line,
  boxrule=0.45pt,
  arc=1.2mm,
  left=3.2mm,
  right=3.2mm,
  top=2.8mm,
  bottom=2.8mm,
  #1
}

\newtcolorbox{softcard}[2][]{
  enhanced,
  colback=#2!10,
  colframe=#2,
  boxrule=0.45pt,
  arc=1.2mm,
  left=3.2mm,
  right=3.2mm,
  top=2.8mm,
  bottom=2.8mm,
  #1
}

\newcommand{\kicker}[1]{%
  {\color{AIEPRed}\scriptsize\bfseries\MakeUppercase{#1}}\par\vspace{2.5mm}
  {\color{AIEPRed}\rule{13mm}{1.2pt}}\vspace{3mm}
}

\newcommand{\labeltext}[1]{%
  {\scriptsize\bfseries\color{AIEPSlate}\MakeUppercase{#1}}\par\vspace{1.3mm}
}

\newcommand{\pagebreaksection}{\newpage}

\begin{document}

% 1 ---------------------------------------------------------------------------
\thispagestyle{fancy}
\kicker{GeoGreen Escolar Osorno}
{\Huge\bfseries\color{AIEPNavy} Taller 2 - Ciencia del reciclaje}\par
\vspace{2mm}
{\large\color{AIEPSlate} Planificación lectiva para facilitar la sesión sobre materiales, reciclabilidad y separación en origen.}\par
\vspace{6mm}

\begin{tabularx}{\textwidth}{Y Y Y}
\begin{softcard}[height=32mm,valign=top]{AIEPSoftBlue}
\labeltext{Duración por bloque}
{\Large\bfseries\color{AIEPNavy}90 minutos}\par
Sesión presencial, participativa y grupal.
\end{softcard}
&
\begin{softcard}[height=32mm,valign=top]{AIEPSoft}
\labeltext{Horario sugerido}
{\Large\bfseries\color{AIEPNavy}A 09:00--10:30}\par
B 10:45--12:15, con 15 minutos de ajuste entre bloques.
\end{softcard}
&
\begin{softcard}[height=32mm,valign=top]{AIEPGreen}
\labeltext{Participantes}
{\Large\bfseries\color{AIEPNavy}2 bloques}\par
30 estudiantes por bloque, organizados en 5 equipos de 6.
\end{softcard}
\end{tabularx}

\vspace{5mm}
\begin{tcolorbox}[
  enhanced,
  colback=AIEPNavy,
  colframe=AIEPNavy,
  boxrule=0pt,
  arc=1mm,
  left=5mm,
  right=5mm,
  top=4mm,
  bottom=4mm
]
{\color{White}\small\bfseries SENTIDO DE LA SESIÓN}\par\vspace{1.5mm}
{\color{White}\large\bfseries El taller transforma la observación del problema ambiental en comprensión del material.}\par
\vspace{1.5mm}
{\color{AIEPSoftBlue}Los equipos aprenden a distinguir objeto y material, reconocer condiciones que permiten o anulan el reciclaje, y construir una ficha que conecte un residuo real con el problema definido en el Taller 1.}
\end{tcolorbox}

\vspace{4mm}
\begin{tabularx}{\textwidth}{Y Y}
\begin{card}[height=35mm,valign=top]
\labeltext{Responsable sugerido}
{\bfseries\color{AIEPNavy}Docente facilitador/a del programa GeoGreen}\par
Conduce la conversación, orienta el trabajo por equipos y cuida que las respuestas se justifiquen desde propiedades del material, no solo desde intuiciones.
\end{card}
&
\begin{card}[height=35mm,valign=top]
\labeltext{Verificadores mínimos}
\begin{itemize}
  \item Lista de asistencia por bloque.
  \item Fotografías del trabajo grupal.
  \item Ficha de material por equipo.
  \item Conexión escrita entre ficha y problema del Taller 1.
\end{itemize}
\end{card}
\end{tabularx}

\pagebreaksection

% 2 ---------------------------------------------------------------------------
\kicker{Propósito formativo}
\section*{Objetivos y aprendizajes esperados}

\begin{softcard}{AIEPSoftBlue}
\labeltext{Objetivo del Taller 2}
Que cada estudiante comprenda que la reciclabilidad de un residuo depende de su composición, de la separación en origen y de las condiciones en que llega al sistema de reciclaje; y que cada equipo pueda caracterizar un material concreto asociado a su problema ambiental.
\end{softcard}

\vspace{3mm}
\begin{tabularx}{\textwidth}{Y Y}
\begin{card}[height=27mm,valign=top]
\labeltext{Resultado central de la sesión}
Cada equipo produce una ficha de material que describe familia, composición, condiciones de reciclaje, factores de contaminación, dudas pendientes y relación con su problema ambiental.
\end{card}
&
\begin{card}[height=27mm,valign=top]
\labeltext{Idea fuerza}
Reciclar bien no empieza en la planta de reciclaje: empieza antes de botar, cuando el residuo se mantiene limpio, seco y separado.
\end{card}
\end{tabularx}

\vspace{4mm}
\subsection*{Resultados de aprendizaje}
\begin{tabularx}{\textwidth}{L{31mm}Y}
\rowcolor{AIEPNavy}
\textcolor{White}{\bfseries Dimensión} & \textcolor{White}{\bfseries Al finalizar el taller, se espera que el/la estudiante pueda...} \\
\hline
Materiales & Distinguir familias de materiales frecuentes: plásticos, papel y cartón, vidrio, metales, orgánicos y envases multicapa. \\
\hline
Reciclabilidad & Explicar por qué un material reciclable puede perderse si llega sucio, mojado o mezclado con otros residuos. \\
\hline
Clasificación & Aplicar criterios simples para decidir si un residuo es reciclable, no reciclable, orgánico, de manejo especial o si genera duda. \\
\hline
Trabajo colaborativo & Completar una ficha de material con acuerdos del equipo, evidencia observable y preguntas pendientes. \\
\hline
Continuidad & Conectar el material estudiado con el problema ambiental definido en el Taller 1 y preparar el puente hacia tecnología en el Taller 3. \\
\hline
\end{tabularx}

\vspace{4mm}
\begin{softcard}{AIEPGold}
\labeltext{Enfoque de evaluación formativa}
No se espera exactitud técnica exhaustiva. El indicador principal es que el equipo avance desde ``esto es basura'' hacia ``esto es un material, se comporta así y se pierde o recupera bajo estas condiciones''.
\end{softcard}

\pagebreaksection

% 3 ---------------------------------------------------------------------------
\kicker{Rol docente}
\section*{Enfoque metodológico y mediación}

\begin{tabularx}{\textwidth}{Y Y}
\begin{card}[height=62mm,valign=top]
\labeltext{Principios de facilitación}
\begin{itemize}
  \item Usar ejemplos concretos y cotidianos.
  \item Preguntar por el material antes que por el contenedor.
  \item Hacer visible la diferencia entre reciclable y reciclado.
  \item Validar la duda cuando falte información.
  \item Cuidar que cada equipo cierre con un producto escrito.
\end{itemize}
\end{card}
&
\begin{softcard}[height=62mm,valign=top]{AIEPGreen}
\labeltext{Rol del/de la docente}
\begin{itemize}
  \item Instalar vocabulario básico sin convertir la sesión en clase de química.
  \item Orientar la observación hacia propiedades y condiciones.
  \item Mediar la clasificación con preguntas de justificación.
  \item Conectar cada ficha con el problema ambiental del equipo.
  \item Preparar el puente hacia medición, alerta y visualización.
\end{itemize}
\end{softcard}
\end{tabularx}

\vspace{4mm}
\begin{softcard}{AIEPRed}
\labeltext{Cuidado pedagógico}
Evitar respuestas absolutas del tipo ``esto siempre se recicla'' o ``esto nunca sirve''. La formulación más útil es condicional: se recupera si está limpio, seco, separado y si existe un canal de gestión adecuado.
\end{softcard}

\vspace{4mm}
\subsection*{Materiales sugeridos}
\begin{tabularx}{\textwidth}{Y Y Y}
\begin{card}[height=35mm,valign=top]
\labeltext{Para activar}
\begin{itemize}
  \item Botella, lata, papel, envase o envoltorio limpio.
  \item Pizarra o papelógrafo.
  \item Plumones.
\end{itemize}
\end{card}
&
\begin{card}[height=35mm,valign=top]
\labeltext{Para trabajar}
\begin{itemize}
  \item Plantilla de ficha de material.
  \item Tarjetas o imágenes de residuos.
  \item Notas adhesivas.
\end{itemize}
\end{card}
&
\begin{card}[height=35mm,valign=top]
\labeltext{Para evidenciar}
\begin{itemize}
  \item Lista de asistencia.
  \item Registro fotográfico.
  \item Fichas completadas por equipo.
\end{itemize}
\end{card}
\end{tabularx}

\vspace{4mm}
\begin{card}
\labeltext{Organización sugerida}
Mantener equipos de 5 a 6 estudiantes con seis responsabilidades persistentes: coordinación, investigación, diseño, tecnología, pruebas y evidencia, y comunicación. Los roles pueden rotar; la propuesta y las decisiones pertenecen al equipo completo.
\end{card}

\pagebreaksection

% 4 ---------------------------------------------------------------------------
\kicker{Secuencia de aula}
\section*{Plan minuto a minuto}

\begin{tabularx}{\textwidth}{L{18mm}L{29mm}Y L{34mm}}
\rowcolor{AIEPNavy}
\textcolor{White}{\bfseries Tiempo} & \textcolor{White}{\bfseries Momento} & \textcolor{White}{\bfseries Actividad docente y estudiantil} & \textcolor{White}{\bfseries Señal de avance} \\
\hline
0--8 min & Encuadre & Presentar el propósito de la sesión: comprender los residuos como materiales y preparar una ficha útil para el proyecto. & Curso reconoce continuidad con Taller 1. \\
\hline
8--20 min & Objeto y material & Abrir con un objeto real. Guiar la distinción entre lo que se usa y el material que define su destino. & Ejemplos de objeto y material en pizarra. \\
\hline
20--35 min & Familias de materiales & Ordenar plásticos, papel/cartón, vidrio, metales, orgánicos y multicapa. Asociar una propiedad clave a cada familia. & Tabla breve construida con aportes del curso. \\
\hline
35--50 min & Reciclable y reciclado & Explicar la cadena de recuperación y la regla limpio, seco y separado. Mostrar cómo la contaminación cambia el destino del residuo. & Estudiantes justifican dos casos. \\
\hline
50--60 min & Colores y sistema & Revisar colores de referencia en Chile, separación en origen, Ley REP y rol de gestores o recicladores de base. & Curso identifica color y condición para materiales comunes. \\
\hline
60--78 min & Ficha de material & Equipos eligen un residuo concreto, completan familia, composición, condiciones, contaminantes y dudas. & Ficha de material avanzada por equipo. \\
\hline
78--86 min & Conexión con problema & Cada equipo relaciona su ficha con la frase del problema del Taller 1. & Conexión escrita entre material y problema. \\
\hline
86--90 min & Cierre & Salida breve: aprendizaje clave, por qué importa entender el material y qué información falta. & Evidencia de cierre y puente al Taller 3. \\
\hline
\end{tabularx}

\vspace{4mm}
\begin{softcard}{AIEPGold}
\labeltext{Manejo del tiempo}
Si el bloque se retrasa, reducir la explicacion del sistema chileno y proteger los últimos 25 minutos. El producto irrenunciable es la ficha de material conectada con el problema del equipo.
\end{softcard}

\pagebreaksection

% 5 ---------------------------------------------------------------------------
\kicker{Desarrollo de actividades}
\section*{Guía práctica para facilitar la sesión}

\begin{tabularx}{\textwidth}{L{25mm}Y}
\begin{softcard}[height=24mm,valign=center]{AIEPSoftBlue}\centering\bfseries Inicio\end{softcard}
&
\begin{card}[height=24mm,valign=top]
\textbf{Instalar la distinción objeto/material.} Usar una botella, lata o envoltorio visible. Preguntar qué ve el curso y luego precisar de qué material está hecho. Cerrar con una idea simple: el objeto es lo que usamos; el material define si puede recuperarse.
\end{card}
\\[2mm]
\begin{softcard}[height=24mm,valign=center]{AIEPSoft}\centering\bfseries Núcleo\end{softcard}
&
\begin{card}[height=24mm,valign=top]
\textbf{Trabajar ``reciclable'' versus ``reciclado''.} Presentar la cadena separar, mantener limpio y seco, entregar, retirar, procesar y convertir en nuevo material. Reforzar que el primer tramo depende de las personas y de los hábitos de separación.
\end{card}
\\[2mm]
\begin{softcard}[height=24mm,valign=center]{AIEPGreen}\centering\bfseries Criterio\end{softcard}
&
\begin{card}[height=24mm,valign=top]
\textbf{Clasificar con justificación.} Usar categorías simples: reciclable, no reciclable, orgánico, manejo especial y genera duda. La categoría ``genera duda'' debe entenderse como una respuesta válida cuando falta información para decidir.
\end{card}
\\[2mm]
\begin{softcard}[height=24mm,valign=center]{AIEPGold}\centering\bfseries Equipos\end{softcard}
&
\begin{card}[height=24mm,valign=top]
\textbf{Construir la ficha de material.} Pedir un residuo concreto, no una categoría amplia. La ficha debe responder: material elegido, familia, composición, condiciones de reciclaje, qué lo arruina, dudas y relación con el problema.
\end{card}
\\[2mm]
\begin{softcard}[height=24mm,valign=center]{AIEPRed}\centering\bfseries Cierre\end{softcard}
&
\begin{card}[height=24mm,valign=top]
\textbf{Conectar material y problema.} Guiar la frase: ``ahora entendemos que el problema no es solo acumulación de residuos; también es pérdida de un material recuperable por mezcla, suciedad, humedad o falta de información''.
\end{card}
\end{tabularx}

\vspace{4mm}
\begin{card}[height=43mm,valign=top]
\labeltext{Preguntas útiles para circular entre equipos}
\begin{tabularx}{\textwidth}{Y Y}
\begin{itemize}
  \item ¿De qué está hecho realmente?
  \item ¿Bajo qué condiciones se recupera?
  \item ¿Qué lo contamina o lo arruina?
\end{itemize}
&
\begin{itemize}
  \item ¿Qué dato falta para decidir mejor?
  \item ¿Cómo se conecta con su problema?
  \item ¿Qué podría medir o avisar una tecnología?
\end{itemize}
\end{tabularx}
\end{card}

\pagebreaksection

% 6 ---------------------------------------------------------------------------
\kicker{Cierre y continuidad}
\section*{Evidencias, criterios de cierre y próximo paso}

\begin{tabularx}{\textwidth}{Y Y}
\begin{card}
\labeltext{Evidencias a recoger}
\begin{itemize}
  \item Lista de asistencia del Bloque A y Bloque B.
  \item Fotografías de activación, trabajo grupal y cierre.
  \item Ficha de material completada por equipo.
  \item Conexión escrita entre ficha y problema.
  \item Dudas o información faltante registradas.
\end{itemize}
\end{card}
&
\begin{softcard}[height=43mm,valign=top]{AIEPGreen}
\labeltext{Criterios de cierre}
\begin{itemize}
  \item La ficha nombra un material concreto.
  \item La familia de material está identificada.
  \item Las condiciones de reciclaje están justificadas.
  \item El equipo reconoce al menos un factor de contaminación.
  \item La conexión con el problema del Taller 1 es comprensible.
\end{itemize}
\end{softcard}
\end{tabularx}

\vspace{4mm}
\begin{softcard}{AIEPSoftBlue}
\labeltext{Salida reflexiva sugerida}
Antes de finalizar, cada estudiante o equipo responde: qué aprendimos sobre nuestro material, por qué importa entenderlo para reciclar bien y qué información falta para decidir con más seguridad.
\end{softcard}

\vspace{4mm}
\begin{tabularx}{\textwidth}{L{27mm}Y}
\rowcolor{AIEPNavy}
\textcolor{White}{\bfseries Continuidad} & \textcolor{White}{\bfseries Cómo se conecta este taller con la secuencia GeoGreen} \\
\hline
Taller 1 & Retoma el problema ambiental formulado por cada equipo y lo mira desde la composición del residuo o material involucrado. \\
\hline
Taller 3 & Prepara la pregunta tecnológica: qué podría medirse, avisarse o visualizarse para mejorar la gestión del problema observado. \\
\hline
Mentorías & Entrega una base para que el equipo desarrolle su propuesta entre sesiones y use las mentorías para revisar, orientar y destrabar avances. \\
\hline
\end{tabularx}

\vspace{4mm}
\begin{tcolorbox}[
  enhanced,
  colback=AIEPNavy,
  colframe=AIEPNavy,
  boxrule=0pt,
  arc=1mm,
  left=4mm,
  right=4mm,
  top=3.5mm,
  bottom=3.5mm
]
{\color{White}\bfseries Mensaje de cierre para la sesión}\par\vspace{1mm}
{\color{AIEPSoftBlue}Entender el material permite mirar el problema con más precisión. Un residuo recuperable puede perderse por mezcla, suciedad o falta de información; una solución sostenible debe responder a esa realidad concreta.}
\end{tcolorbox}

\end{document}

